\section{Motivación y antecedentes}

En el actual plan de estudios del Grado de Ingeniería Informáica en la asignatura de Internet de las cosas se estudia el ecosistema de los sistemas empotrados y sus múltiples protocolos de comunicación. Dentro de la amplia gama de protocolos de comunicación el estudio del bus CAN y se comentó su uso en el sector automotriz, esto dió lugar a la motivación por la cual se ha desarrollado este Trabajo de Fin de Título en el que se estudia el uso del bus CAN y el protoclo OBD-II y demás protocolos que orquestan la comunicación de información en un vehiculo.

Actualmente los automóviles cuentan con un sistema de numerosos microcontroladores que generan una cantidad muy voluminosa de datos de telemetría y además cuentan con el protocolo de diagnóstico abordo (OBD-II). Estos son accesbles a traves del puerto OBD que se encuentra normalmente en la parte inferior del volante en el asiento del piloto. Este puerto permite que no solo los talleres y personas especializadas en el mundo de la mecánica y electromecánica, sino que cualquier persona con un "escáner" pueda ver el estado de su coche o pueda diagnositcar que fallo tiene o anticipar posibles fallos.

Actualmente en el mercado han aparecido numerosas soluciones para poner de forma accesible el acceso al diagnóstico OBD-II. La mayor parte de estas soluciones usan el adaptador ELM327, montando un circuito integrado que actúa como puente entre el bus CAN del vehículo y una interfaz accesible vía Bluetooth o Wi-Fi. Por otra parte existen numerosas aplicaciones móviles compatibles con estos adaptadores tales como Torque Pro u OBD Fusion, no obstantes estas soluciones no suelen permitir el registro estructurado de datos para análisis y la mayor parte de estas soluciones introducen latencias significativas dano lugar a una experiencia de usuario pobre.

\section{Objetivos}
\label{sec:objetivos}

A partir del contexto descrito, se definen los siguientes objetivos para este proyecto.

\subsection{Objetivo general}

Diseñar y desarrollar un sistema empotrado de diagnóstico del vehículo basado en el protocolo OBD-II sobre bus CAN, con acceso inalámbrico vía Bluetooth Low Energy, de bajo coste económico y validado mediante un banco de pruebas basado en Arduino.

\subsection{Objetivos específicos}

La propuesta de este trabajo (TFT01) define tres funciones de usuario de alto nivel que el sistema debe satisfacer. La tabla~\ref{tab:tft01_objetivos} recoge estas funciones junto con los objetivos específicos de la memoria que las implementan.

\begin{table}[!htbp]
\caption{Trazabilidad entre funciones definidas en la propuesta TFT01 y objetivos específicos del trabajo.}
\label{tab:tft01_objetivos}
\centering
\renewcommand{\arraystretch}{1.3}
\begin{tabular}{|p{4.2cm}|p{7.5cm}|c|}
\hline
\textbf{Función TFT01} & \textbf{Descripción} & \textbf{Obj.} \\
\hline
Monitorización en tiempo real &
  Visualización numérica de parámetros de sensores disponibles en el vehículo (RPM, temperaturas, presiones, voltajes, etc.) &
  O1, O2, O4, O5, O6 \\
\hline
Consulta y gestión de errores &
  Lectura y borrado de códigos de error (DTCs) almacenados en la memoria de las ECUs &
  O1, O2, O3, O4, O6 \\
\hline
Identificación &
  Consulta de datos del vehículo (VIN) e información de los distintos módulos integrados &
  O1, O2, O4, O5 \\
\hline
\end{tabular}
\end{table}

Los seis objetivos específicos que concretan estas funciones son:

\begin{enumerate}
    \item[\textbf{O1.}] \textbf{Implementar la pila de protocolos OBD-II completa}: desarrollar sobre Raspberry Pi una implementación funcional de ISO-TP (ISO 15765-2) como capa de transporte y OBD-II (SAE J1979) como capa de aplicación, cubriendo los modos de servicio 0x01 (datos en vivo), 0x03 (lectura de DTCs), 0x04 (borrado de DTCs) y 0x09 (información del vehículo, incluyendo VIN).

    \item[\textbf{O2.}] \textbf{Desarrollar un simulador de ECU}: implementar sobre Arduino MKR con shield CAN un simulador de ECU que responda a peticiones OBD-II estándar con datos de telemetría realistas, con el objetivo de disponer de un banco de pruebas reproducible y controlado que no dependa de la disponibilidad de un vehículo real.

    \item[\textbf{O3.}] \textbf{Diseñar mecanismos de simulación de condiciones adversas}: incorporar en el simulador Arduino la generación de ruido CAN (pérdida de tramas, latencia variable, respuestas negativas esporádicas) y tráfico de fondo aleatorio, para validar la robustez de la herramienta de diagnóstico.

    \item[\textbf{O4.}] \textbf{Proporcionar acceso inalámbrico mediante BLE}: implementar un servidor BLE GATT sobre la Raspberry Pi que exponga la funcionalidad de diagnóstico a través del servicio Nordic UART Service (NUS), eliminando la necesidad de conexión física entre el dispositivo de diagnóstico y el usuario.

    \item[\textbf{O5.}] \textbf{Desarrollar una aplicación móvil multiplataforma}: implementar una aplicación para Android e iOS con React Native que permita al usuario conectarse al sistema de diagnóstico vía BLE, monitorizar parámetros del vehículo en tiempo real, gestionar los códigos de avería y consultar el historial de sesiones almacenadas.

    \item[\textbf{O6.}] \textbf{Registrar y almacenar los datos de diagnóstico}: implementar un sistema de logging estructurado en base de datos SQLite que almacene las sesiones, las muestras de telemetría y los comandos ejecutados, para su consulta y análisis posterior.
\end{enumerate}

\subsection{Alcance y limitaciones}
\label{sec:alcance}

El sistema desarrollado se centra en la implementación de los modos OBD-II estándar definidos por SAE J1979, sin abordar los servicios de diagnóstico propietarios de cada fabricante, que requieren herramientas específicas y, en ka mayor parte casos, acuerdos de licencia con el fabricante del vehículo. Asimismo, la validación experimental del sistema se realiza sobre el banco de pruebas desarrollado en Arduino, dado que las restricciones de acceso al hardware del vehículo real durante el período de desarrollo del trabajo no permitieron completar las pruebas sobre el vehículo objetivo. Las implicaciones de esta limitación y las líneas de trabajo futuro que abre se discuten en el Capítulo~7.
